\chapter{Einleitung}
\label{cha:Einleitung}

Wenn man automatisiert handeln möchste braucht man eine Stragie. Zum testen braucht man auch etwas. Den Backtester.
Arbeit wird ein Backtster entwickelt. Der clientseitig ist. Weniger Sicherheitslücken. Doch wie ist der Backtester am schnellsten.
Deswegen rechenintenveste Teil ausgelagert. Die Strategien. Diese mit neuen Web Technologien wie WASM im Vergleich zu reinem JS.
Nur diese 1 Seite

\section{Motivation der Arbeit}
\label{sec:MotivationderArbeit}

\section{Ziele der Arbeit}
\label{sec:ZielederArbeit}

\section{Vorgangsweise zur Zielerreichung}
\label{sec:VorgangsweisezurZielerreichung}

\section{Kapitelübersicht}
\label{sec:Kapitelübersicht}

In diesem Unterkapitel werden die einzelnen Kapitel überblicksmäßig beschrieben.

\subsection{Grundlegende Konzepte und Technologien}
\label{subsec:GrundlegendeKonzepteundTechnologien}

In dem Kaptiel \ref{cha:GrundlegendeKonzepteundTechnologien} werden die Grundlagen und die relevanten Technologien beschrieben.
Dabei werden die grundlegenden Themen wie Backtesting, Biase im Backtesting, bereits existierende Backtester sowie Ansätze für Handelsstrategien näher erläutert.
Zudem wird WebAssembly als eine der Kerntechnologien vorgestellt und wie diese in Webanwendungen integriert wird.

\subsection{Konzept und Design}
\label{subsec:KonzeptundDesign}

Das Kaptiel \ref{cha:KonzeptundDesign} beschreibt wie die einzelnen Komponenten des entwickelten Backtester aufgebaut sind, wie die Geschwindigkeitstests durchgeführt werden und welche Werkzeuge für die Entwicklung von WebAssembly genutzt wurden. 

\subsection{Implementierung des Backtesters und der Strategien}
\label{subsec:ImplementierungdesBacktestersundderStrategien}

In dem Kaptiel \ref{cha:ImplementierungdesBacktestersundderStrategien} wird die Implementierung des Backtesters und der Handelsstrategien beschrieben.
Dabei wird der gesamte Prozess von der Beschaffung historischer Daten über das Laden und Anwenden der Handelsstrategien bis hin zur Darstellung der Ergebnisse auf der Webseite beschrieben.

\subsection{Performanztests der unterschiedlichen Programmiersprachen}
\label{subsec:PerformanztestsderunterschiedlichenProgrammiersprachen}

Das Kaptiel \ref{cha:PerformanztestsderunterschiedlichenProgrammiersprachen} zeigt den Aufbau und die Ergebnisse der Performanztests der verschiedenen Implementierungen der Handelsstrategien in unterschiedlichen Programmiersprachen.
Es werden Strategien in verschiedene Programmiersprachen implementiert und deren Ausführungszeiten im Backtester in verschiedenen Browsern gemessen.
Desweitern werden die Ergebnisse der Performanztests analysiert und interpretiert.

\subsection{Zusammenfassung und Schlussbemerkungen}
\label{subsec:Zusammenfassung und Schlussbemerkungen}

Im Kapitel \ref{cha:ZusammenfassungundSchlussbemerkungen} wird die Arbeit zusammengefasst, ein Ausblick auf mögliche zukünftige Arbeiten gegeben und abschließend Erfahrungen einbezogen.
